\paragraph{Algebraic Rigidity of the SM Signature, Adjacent-Layer Activation Prohibition, and Carry Absorption}\label{par:bdry-tower-sm-rigidity-adjacent-carry}
This paragraph promotes the \(D=12\) counting interface of the boundary three-layer tower (\S\ref{subsec:bdry-tower-zeck-gut}) together with the Zeckendorf--Fibonacci signature dictionary (Table \ref{tab:zeckendorf_gut_signatures}) into three citable structural assertions: one concerning the compact Lie algebra rigidity of \(D=12\), one concerning the impossibility of strict inclusion due to adjacent-layer activation prohibition, and one concerning the carry absorption of the \(m=9\) excitation at the audit signature level.

\begin{proposition}[SM Gauge Algebra Rigidity under Zeckendorf Atomicity]\label{prop:sm-zeckendorf-lie-algebra-rigidity}
Under the caliber where ``the total number of gauge bosons \(D\) is given by boundary tower layer activation'' and ``no redundant activation'' is audited via Zeckendorf unique decomposition,
if \(D=12\), then its Zeckendorf decomposition is uniquely
\[
12=F_{6}+F_{4}+F_{2},
\qquad
\Sigma(\mathfrak g_{\mathrm{SM}})=\{8,6,4\}
\]
(Table \ref{tab:zeckendorf_gut_signatures}).
Furthermore, if one requires ``term-by-term alignment of each factor dimension with Zeckendorf terms'' (i.e., each activated layer corresponds to an irreducible compact simple/abelian factor dimension slot),
then any compact Lie algebra \(\mathfrak g\) compatible with this audit must be isomorphic to
\[
\boxed{\ \mathfrak g\ \cong\ \mathfrak u(1)\ \oplus\ \mathfrak{su}(2)\ \oplus\ \mathfrak{su}(3)\ }.
\]
Moreover, this decomposition is unique in the sense of ``term-by-term alignment.''
\end{proposition}
\begin{proof}
By Zeckendorf uniqueness, the decomposition of \(12\) can only be \(F_6+F_4+F_2=8+3+1\), with repeated and adjacent terms prohibited.
In the classification of compact Lie algebras, the compact Lie algebra of dimension \(1\) must be \(\mathfrak u(1)\);
the compact simple Lie algebra of dimension \(3\) is unique up to isomorphism (\(\mathfrak{su}(2)\cong\mathfrak{so}(3)\));
and the compact simple Lie algebra of dimension \(8\) is also unique (\(\mathfrak{su}(3)\)).
Under the ``term-by-term alignment'' caliber, the three dimension slots can only be filled by the above three factors, yielding the stated direct sum decomposition, and uniqueness follows. \qed
\end{proof}

\begin{corollary}[Adjacent-Layer Activation Prohibition and Impossibility of Strict Inclusion]\label{cor:zeckendorf-no-adjacent-strict-inclusion-impossible}
For any candidate compact Lie algebra \(\mathfrak g\), its resolution signature \(\Sigma(\mathfrak g)=\{k+2:k\in S\}\) is given by the Zeckendorf index set \(S\), which contains no adjacent indices (\S\ref{subsec:bdry-tower-zeck-gut}).
Therefore for all \(m\),
\[
m\in\Sigma(\mathfrak g)\ \Longrightarrow\ m\pm 1\notin\Sigma(\mathfrak g).
\]
In particular, under the strict caliber requiring
\(
\Sigma(\mathfrak g)\supseteq \Sigma(\mathfrak g_{\mathrm{SM}})=\{4,6,8\}
\)
(Table \ref{tab:zeckendorf_gut_signatures}), it necessarily follows that
\(
\{5,7\}\cap \Sigma(\mathfrak g)=\varnothing
\).
Hence any candidate whose signature contains \(m=5\) or \(m=7\) cannot strictly include SM at the signature level; for example, the signature of \(E_6\) is \(\Sigma(E_6)=\{12,10,5\}\) (Table \ref{tab:zeckendorf_gut_signatures}), which does not satisfy strict inclusion.
\end{corollary}
\begin{proof}
By the Zeckendorf condition ``no adjacent indices in the index set,'' if \(m=k+2\in\Sigma(\mathfrak g)\), then \(k\pm 1\notin S\), equivalently \(m\pm 1\notin \Sigma(\mathfrak g)\).
Applying this to \(\Sigma(\mathfrak g)\supseteq\{4,6,8\}\) immediately yields \(\{5,7\}\cap\Sigma(\mathfrak g)=\varnothing\); the last statement follows by direct verification from Table \ref{tab:zeckendorf_gut_signatures}. \qed
\end{proof}

\begin{corollary}[Separation of Protocol Excitation Layer and Audit Signature Layer: Carry Absorption at \texorpdfstring{$m=9$}{m=9}]\label{cor:zeckendorf-carry-absorption-m9}
With the SM background signature \(\Sigma(\mathfrak g_{\mathrm{SM}})=\{4,6,8\}\) already fixed,
any construction that attempts to additionally activate \(m=9\) (corresponding to \(F_7=13\) additional boundary modes) cannot be maintained as a ``non-redundant new layer'' at the Zeckendorf audit signature level:
once the layer set is normalized to Zeckendorf normal form, \(m=8\) and \(m=9\) trigger the unique carry absorption
\[
F_{7}+F_{6}=F_{8},
\qquad\text{equivalently}\qquad
\{8,9\}\ \leadsto\ \{10\}.
\]
Therefore, \(m=9\) can only exist as a transient protocol-layer excitation; in the minimal non-redundant audit sense of \(\Sigma(\cdot)\), the minimal upward jump must land at \(m=10\).
\end{corollary}
\begin{proof}
The definition of \(\Sigma(\mathfrak g)\) requires the Zeckendorf index set to contain no adjacent indices (\S\ref{subsec:bdry-tower-zeck-gut}).
Since the SM background contains \(m=8\) (i.e., contains the \(F_6\) term), adding \(m=9\) (i.e., adding the \(F_7\) term) produces adjacent indices \(6,7\), violating the normal form condition.
The unique Zeckendorf normalization for this adjacent pair is \(F_6+F_7=F_8\), corresponding to the layer replacement \(\{8,9\}\leadsto\{10\}\).
Hence \(m=9\) does not form an independent new layer at the audit signature level, and the conclusion holds. \qed
\end{proof}

\begin{corollary}[Resonance Ladder -- Four-Step Shift -- Minimal Triple Selection Law under Zeckendorf Constraint]\label{cor:sm-minimal-triple-selection-law}
Under the boundary word tower caliber, if one requires ``at least two non-abelian resonance layers'' to appear in the resolution signature as
\(
\{6,8\}\subseteq \Sigma(\mathfrak g)
\)
(see Corollary \ref{cor:fib-lie-resonance-scarcity-su2-su3} and Remark \ref{rem:fib-lie-resonance-shift-by-4} for the \(m\mapsto m+4\) shift correspondence), and the layer set is taken in Zeckendorf normal form with ``no adjacent indices, no redundant activation,''
then the arithmetically minimal solution for \(\Sigma(\mathfrak g)\) is unique:
\[
\Sigma(\mathfrak g)=\Sigma(\mathfrak g_{\mathrm{SM}})=\{4,6,8\},
\qquad
D=\sum_{m\in\Sigma(\mathfrak g)}|X_m^{\mathrm{bdry}}|=1+3+8=12.
\]
Hence under the term-by-term alignment caliber, the isomorphism type of \(\mathfrak g\) is forced to be
\(
\mathfrak u(1)\oplus\mathfrak{su}(2)\oplus\mathfrak{su}(3)
\)
(Proposition \ref{prop:sm-zeckendorf-lie-algebra-rigidity}).
\end{corollary}
\begin{proof}
By Corollary \ref{cor:fib-lie-resonance-scarcity-su2-su3}, the stable type resonance ladder at the minimal window gives the earliest non-abelian resonance points for \(\dim\mathfrak{su}(2)=3\) and \(\dim\mathfrak{su}(3)=8\);
by Remark \ref{rem:fib-lie-resonance-shift-by-4}, under the boundary tower caliber these two points correspond to boundary layers \(m=6,8\), so ``at least two non-abelian resonance layers'' forces \(\{6,8\}\subseteq\Sigma(\mathfrak g)\).
Zeckendorf normal form prohibits adjacent layers from coexisting, hence \(m=7,9\) cannot coexist with \(\{6,8\}\).
Among the remaining admissible layers, the smallest is \(m=4\) (contributing \(|X_4^{\mathrm{bdry}}|=F_2=1\)), yielding the unique minimal signature \(\{4,6,8\}\);
by \(|X_m^{\mathrm{bdry}}|=F_{m-2}\) one obtains \(D=1+3+8=12\).
Finally, by Proposition \ref{prop:sm-zeckendorf-lie-algebra-rigidity}, the isomorphism type and uniqueness of \(\mathfrak g\) follow. \qed
\end{proof}
